\chapter{Contextualizing the Napatan period: historical background and cultural interactions}

In this chapter, I will examine the historical context of the Napatan period, as well as the history of interactions between Nubia and Egypt. I begin with an outline of the period, tracing its origins from the withdrawal of Egyptian administration at the end of the New Kingdom and to the last known rulers of the Napatan era. Both the beginning and end of the period are characterized by a scarcity of material evidence. Notably, for the latter part of the period, we know little beyond the names of its rulers, leaving much about the region's overall history still unknown.

Within this context, it is necessary to explore the origins and early chronology of the Napatan dynasty—a question in the center of scholarly debate for many years. The various theories surrounding this topic are important to understanding the processes of "Egyptianization" and cultural entanglements in the Nile Valley region. Specifically, these discussions focus on whether the emergence of the dynasty was an internal development or the result of external, i.e. Egyptian, influences.

Following this discussion and to set the stage for a detailed examination of amulets during the Napatan period, I will outline the complex and historical relations between Nubia and Egypt, spanning from the 4th millennium BCE to the Napatan period (ca. 750—270 BCE). These interactions are particularly important for grasping why the concept of cultural entanglements offers a more nuanced framework for interpreting not only the presence of “Egyptian-style” features and material culture in Nubia but also the persistence of Nubian objects and motifs, and even the synthesis of both into what can be termed a "third space” (See Chapter 2). Moreover, by focusing on recent archaeological discoveries and material evidence, we can gain insights into the experiences of non-elite social groups that have often been overlooked in traditional written and visual sources.

Funerary amulets were also part of this complex interplay, and by examining the prolonged cultural interactions between Nubia and Egypt over many centuries and across diverse cultures, we will be better equipped to analyze and discuss the use of amulets in Napatan tombs across different social groups later on.

\section{Nubia during the Napatan Period}

The Napatan period can be characterized as a shift in Kush’s position in the Nile Valley during the first half of the first millennium BCE. After almost five hundred years of being under Egyptian control during the New Kingdom (ca. 1550—1070 BCE), a new dynasty emerged during the 8th century BCE. They not only conquered Upper Egypt, but also Lower Egypt, unifying the territory as the 25th Dynasty (ca. 715—664 BCE) during the Third Intermediate Period when different rulers controlled each region.

The period is named after the burial sites of its kings, El-Kurru and Nuri, located near the ancient city of Napata \parencite[p.~1]{popeNapatanPeriod2020}. The origin and precise chronology of the Napatan period has been the subject of different theories and debates for many years. After the Egyptian withdrawal from Lower Nubia ca. 1070 BCE, material and textual evidence is extremely scarce in the region for the next three centuries. For many years, scholars have referred to the period following the New Kingdom as the 'dark ages' of Nubia, due to the scarcity of archaeological evidence \parencites[p.~250]{adamsNubiaCorridorAfrica1977}[pp.~88, 109-110]{torokKingdomKushHandbook1997}[p.~129]{morkotBlackPharaohsEgypts2000}[pp.~412-413]{williamsNapatanNeoKushiteState2020}. Earlier authors even proposed that the temple towns of Lower Nubia were depopulated since the Egyptians had returned to their native land while Nubians withdrew to Upper Nubia \parencites[p.~108-109]{adamsPostPharaonicNubiaLight1964}{firthArchaeologicalReportNubia1912}[p.109]{edwardsNubianArchaeologySudan2004}. It has been suggested that this was likely a misinterpretation of the archaeological material from post New Kingdom, which was being identified as either New Kingdom or Napatan \parencites[p.~108-109]{adamsPostPharaonicNubiaLight1964}{firthArchaeologicalReportNubia1912}[p.~109]{edwardsNubianArchaeologySudan2004}.

It has been suggested that this was likely a misinterpretation of the archaeological material from post New Kingdom, which was being identified as either New Kingdom or Napatan \parencites[p.~109-110]{torokKingdomKushHandbook1997}[p.~37]{williamsExcavationsAbuSimbel1990}[p.~415]{williamsNapatanNeoKushiteState2020}. Indeed, recent archaeological work indicate that some cemeteries were occupied continuously from the New Kingdom to the Napatan period, such as Tombos \parencites{buzonEntanglementFormationAncient2016}[p.~7]{smithEthnicityCulture2007}, Hillat el-Arab \parencite{vincentelliHillatElArabJoint2006}, Qustul, Adindan \parencite[p.141]{williamsExcavationsAbuSimbel1992}, Debeira 176 \parencite[p.~200-205]{save-soderberghScandinavianJointExpedition1989}, Amata West \parencites[p.~57]{spencerCemeteriesLateRamesside2009}{binderCemeteryAmaraWest2011}{binder10th9thCenturyBC2011}, and Sanam \parencite{lohwasserAspekteNapatanischenGesellschaft2012}. Settlements and fortresses have also been dated to the post New Kingdom period \parencites[p.~1071-1078]{jesseFendingDesertDwellers2019}[p.~416-417]{williamsNapatanNeoKushiteState2020}.

Despite the limited evidence available, Lower Nubia likely remained populated after the Egyptian withdrawal. Given their role as administrators under Egyptian rule, it is plausible that local elites continued to govern independently and contributed to the rise of the Napatan state \parencites[p.~116]{edwardsNubianArchaeologySudan2004}[p.~111]{torokKingdomKushHandbook1997}. A textual source was found in the Egyptian fortress of Semna. In the inscription, an anonymous king relates his problems with rebel chiefs to a queen, Katimala, seeking her help \parencites{darnellInscriptionQueenKatimala2006}[p.~296]{torokTwoWorldsFrontier2009}. The text has been dated palaeographically to the 21st or 22nd Dynasties (ca. 1069—720 BCE). It has been suggested that the content implies that parallel Nubian kingdoms had developed in the region during the post New Kingdom period. John C. Darnell, who dated the inscription, suggests it may relate to the emergence of the Napatan dynasty \parencite[p.~60]{darnellInscriptionQueenKatimala2006}. However, other scholars have criticized the dating, proposing a date closer to the 25th Dynasty (ca. 744—656 BCE) \parencite[p.~298]{torokTwoWorldsFrontier2009}. At the current state of knowledge, Queen Katimala’s inscription may or may not be a clue as to the political organization of Nubia during the post New Kingdom period and the eventual rise of the Napatan kings.

Nonetheless, by the 8th century BCE, a unified kingdom had been established in Upper Nubia which would eventually expand north and reach Lower Egypt. During the 25th Dynasty, textual evidence written in Egyptian hieroglyphic script and language along with Manetho’s list of kings, provided scholars with written records of the names of the Kushite rulers \parencites{payraudeauRetourSuccessionShabaqoShabataqo2014}{pope25thDynasty2019}.\footnote{Scholars have proposed different absolute dates for the 25th Dynasty rulers \parencites{depuydtDatePiyesEgyptian1993}{kahnInscriptionSargonII2001}{hornungKushiteRulersPre25th2006}{payraudeauRetourSuccessionShabaqoShabataqo2014}. For this thesis, we will use Payraudeau’s suggested chronology.}

The first documented Kushite ruler is Alara, whose name only appears in inscriptions of his successors \parencite[p.~110]{emberlingKushDynastyNapata2023}. Alara was succeeded by Kashta\footnote{Although the name Kashta has been interpreted to mean “the Kushite”, \textcite{vinogradovDidNameKashta2003} demonstrates that the words “Kush” and “Kashta” follow distinct spelling logic.}, who installed his daughter, Amenirdis I, as God’s Wife of Amun at Thebes, establishing his control over Upper Egypt \parencites[p.~319]{torokTwoWorldsFrontier2009}[p.~110]{emberlingKushDynastyNapata2023}.\footnote{On God’s Wives of Amun during the 25th Dynasty, see \cite{ayadGenderRitualManipulation2016}, \cite{beckerFemaleInfluenceAside2016}, \cite{lohwasserNubianessGodsWives2016}.} 
His successor, Piankhy (ca. 744—714 BCE), adopted the traditional Egyptian five-part royal titles \parencites[p.~153]{torokKingdomKushHandbook1997}[p.~324]{torokTwoWorldsFrontier2009}.\footnote{The name Piye is also commonly used in the literature; however, \textcite{rillyNouvelleInterpretationNom2001} argues in favour of using Piankhy.} 
In his Sandstone Stela, likely erected in Year 3 of his reign, Piankhy declared himself the legitimate ruler of Egypt and intended to keep Egypt’s political map divided as long as the other rulers recognized his supremacy and payed him tribute \parencite[p.~155]{torokKingdomKushHandbook1997}.The following kings, Shabataqo and Shabaqo\footnote{It was previously believed that Shabaqo was followed in succession by Shabataqo. Recent scholarship suggests reversing the succession order, supported by Kushite and Egyptian textual evidence \parencites[p.~4-9]{payraudeauRetourSuccessionShabaqoShabataqo2014}{jurmanOrderKushiteKings2017}. Analysis of horse burials at El-Kurru further reinforces this view, citing inscribed cartouches with the kings’ names \parencite{naserKingsHorsesTwo2020}.}, relocated the royal court to Memphis. Shabaqo further consolidated the double kingdom by controlling the local rulers of the Delta region \parencites[p.~166-167]{torokKingdomKushHandbook1997}[p.~141]{rillyHistoireSoudanOrigines2017}.\footnote{A special case is Tanis, where no local rulers were documented and was ruled directly by Shabaqo \parencite{mojeHerrschaftsraumeUndHerrschaftswissen2013}.}
Taharqo succeeded Shabaqo and ruled the double kingdom, building and enlarging temples in Egypt and Kush (ca. 690—664 BCE).\footnote{For a more detailed study on Taharqo’s double kingdom, see \cite{dalliborTaharqoPharaoAus2005}, \cite{popeDoubleKingdomTaharqo2014}.}. However, three years after successfully driving out Assyrian troops in Palestine, Memphis was invaded and sacked. Taharqo, wounded, fled to Thebes. The Assyrians installed new rulers—both Assyrian and Egyptian—in Lower Egypt, effectively controlling the whole region. But Taharqo succeeded in regaining control of Memphis. This victory was, however, short-lived. The Assyrians, ca. 667 BCE, launched a campaign that not only took back control of Lower Egypt, but also defeated the Kushite-Egyptian forces in Upper Egypt. Taharqo, this time, fled to Napata where he died and was the first king buried at Nuri \parencite[p.~145-154]{rillyHistoireSoudanOrigines2017}. His successor, Tanwetamani, claims to have briefly regained control of Lower and Upper Egypt, before being driven away to Kush when the Assyrians not only conquered Upper Egypt again, but looted and burned Thebes. The double kingdom was effectively over, even though commercial and trade relations were maintained between Egypt and Kush \parencite[p.~356-360]{torokTwoWorldsFrontier2009}. Tanwetamani was buried at El-Kurru and  not in the new established cemetery at Nuri. Almost all succeeding kings were buried at Nuri, with the exception for Taharqa’s sucessor, Tanwetamani, who still built his tomb (Ku. 16)\footnote{Only one tomb is thought to have been built at El-Kurru again, pyramid Ku. 1. Since the tomb had no inscribed material, it is not known who was the supposed owner. Recent excavations agree that it was build in the late Napatan or early Meroitic periods and was left unfinished. The excavators tentatively interpret that this pyramid remained unused due to political conflicts among rulers in Napata. These conflicts could have led to the transition from Nuri as the royal burial site to Meroe at the end of the 4th century BCE \parencites[p.~62]{emberlingExcavationPyramidKu2015}{yellinPyramidChapelDecorations2015}.}

The following four kings, Atlanersa, Senkamanisken, Anlamani, and Aspelta, are considered the first rulers of the “Napatan kingdom” and collectively reigned for approximately eighty years. During their reign, they built and expanded temples, commissioned statues of themselves and earlier rulers such as Taharqo and Tanwetamani, and issued public inscriptions and seals \parencites[p.~157-163]{rillyHistoireSoudanOrigines2017}[p.~133-137]{emberlingKushDynastyNapata2023}. Around 593 BCE, Psamtek II, the king of Egypt, invaded Kush. Most scholars believe this invasion occurred during Aspelta's reign, as the sequence of destroyed Kushite royal statues ends with him. The burning of the “Treasury” at Sanam is also considered a consequence of this raid \parencites[p.~361-362]{torokTwoWorldsFrontier2009}{andersonWhatAreThese2009}[p.~138-139]{emberlingKushDynastyNapata2023}.

Apart from their names and tombs at Nuri, very little is known about the approximately seventeen succeeding kings \parencites[p.~435]{williamsNapatanNeoKushiteState2020}[p.~139]{emberlingKushDynastyNapata2023}.\footnote{For the complete list see \cite[p.~496-497]{hornungKushiteRulersPre25th2006}.} They are followed by a group of five rulers known as “Neo-Ramesside kings” due to their titles and inscriptions bearing similarities to the Egyptian Ramesses. The dating of these rulers has been a subject of debate. Some scholars have placed them as the last kings of the Napatan period due to the paleography and contents of their titles \parencite[p.~293]{torokTwoWorldsFrontier2009}. Others have argued that they should be placed before the 25th Dynasty, since the Libyan kings of the 23rd Dynasty also emulated Ramesside titles \parencites[p.~216-217]{morkotEmptyYearsNubian1991}[p.~145-150]{morkotBlackPharaohsEgypts2000}[p.~63-64]{kendallOriginNapatanState1999}. Recently, in an article about pre-25th Dynasty kings, Angelika Lohwasser published a comprehensive reassessment of the textual evidence concerning the Neo-Ramesside kings, offering new insights into their historical context. Lohwasser tentatively concludes that only two of these kings could possibly be attributed to the end of the Napatan period \parencite[p.~132-134]{lohwasserKushiteKingsTwentyfifth2022}. Therefore, the dating of these kings remains uncertain due to the limited and inconclusive nature of the available evidence.

Regardless of who were the last kings of the period, the primary event signaling the end of the Napatan dynasty was the relocation of the royal cemetery from Nuri to Meroe with the burial of Arkamani I around 270 BCE. This change also shifted the political and economic center from Napata to Meroe. This transition to the Meroitic period—which lasted until ca. 350 AD—did not involve abrupt cultural changes, but new cultural developments that were gradually emerging such as the cult of local deities, ruling queens known as kandake, and the  introduction of the Meroitic script \parencites[p.~144]{elhassanRoleEarlyNapatan2023}[p.~2]{kuckertzMeroeEgypt2021}.

\section{Tracing the Napatan Dynasty: chronological and emergence issues}

The origins of the Napatan rulers has been the subject of debate for many years.\footnote{For a comprehensive overview of the literature on the many theories regarding the Napatan dynasty origins and initial chronology, see \cite{morkotPriestlyOriginNapatan2003}.} In the centre of the discussions is the cemetery of El-Kurru, which is considered as the burial site of the ancestors of the Napatan rulers \parencite{torokKingdomKushHandbook1997, torokTwoWorldsFrontier2009, morkotKingshipKinshipEmpire1999}. The importance of El-Kurru is twofold. The first is the emergence question and the second is the overall chronology of the period and, consequently, ruling dates of the kings. As Török notes, “The different el-Kurru chronologies describe and at the same time create different images of political and religious processes” \parencite[p.~304]{torokTwoWorldsFrontier2009}. Scholars have suggested different models for the emergence of the Napatan dynasty and chronology of el-Kurru—known as the short and long chronology.

The history of the cemetery of El-Kurru was first established by its excavator George A. Reisner and published by Dows Dunham in 1950. In the El-Kurru volume of his series “The Royal Cemeteries of Kush”, Dunham published a chronology based on Reisner’s with some of his own modifications, eliminating some original inconsistencies of both Napatan and Meroitic periods, dividing the rulers by generations \parencite[p.~2]{dunhamElKurru1950}. This division was adopted by Dunham in all of his publications on Kushite cemeteries. At El-Kurru, Reisner identified sixteen tombs and divided them into six different evolutionary types. Starting at ca. 860—840 BCE, the tombs would, therefore, be of six different generations of rulers and their wives \cite{reisnerDiscoveryTombsEgyptian1919, reisnerNoteHarvardBostonExcavations1920}. This is known as the “short chronology” and was followed by Timothy Kendall. Kendall, however, proposed the date of ca. 885—835 BCE for the first tombs and made some modifications in the order of the burials according to a tomb evolution based on new evidence \parencite[p.~4]{kendallOriginNapatanState1999}. Regarding the origin of the rulers, Kendall argues that the Egyptian influences and cultural markers at El-Kurru correspond to the influx of priests that rebelled against the new high priest of Amun in ca. 839 BCE and fled to the south \parencite[p.~5,57]{kendallOriginNapatanState1999}. Therefore, according to Kendall, the El-Kurru ancestral rulers were a non-Egyptianized Nubian group that, through the external influence of native Egyptians, were rapidly “Egyptianized”.

Kendall’s model was criticized by Lazló Török, who proposed a longer chronology \parencite[p.~298-309]{torokEmergenceKingdomKush1995, torokOriginNapatanState1999, torokTwoWorldsFrontier2009}. He argues for an earlier date for the first burials, around 995—975 BCE, and asserts that the central sixteen burials belonged to ruling males \parencite[p.~301]{torokTwoWorldsFrontier2009}. Regarding the emergence of the Napatan rulers, Török suggests that the Nubian elite, previously integrated into the New Kingdom Egyptian administration, maintained their power after the withdrawal, taking advantage of the direct access to gold-mining sites \parencite[p.~111-112]{torokKingdomKushHandbook1997}. He characterizes the Napatan state as a transformation from chiefdom to kingdom, with its Egyptianized features resulting from the fusion of indigenous and Egyptian concepts \parencites{torokKingdomKushHandbook1997, torokTwoWorldsFrontier2009}. Other scholars have also supported this theory on the origins of the Napatan dynasty \parencites{morkotFoundationsKushiteState1995a, morkotEgyptNubiaEgyptian2001}{yellinEgyptianReligionIts1995}{smithNubiaEgyptInteraction1998}{binder10th9thCenturyBC2011}. Table \ref{tab:el-kurru-tombs} presents both short and long chronologies of the ancestral El-Kurru burials alongside their proposed occupants.

{\small
	\begin{xltabular}{\textwidth}{X X X X X X}
		\caption{El-Kurru ancestral tombs with their respective dating and occupant according to the short and long chronology}
		\label{tab:el-kurru-tombs} \\
		\toprule
		\multicolumn{3}{c}{\textbf{Kendall's short chronology (1999)}} & \multicolumn{3}{c}{\textbf{Török's long chronology (2009)}} \\
		\cmidrule(r){1-3} \cmidrule(l){4-6}
		\textbf{Tomb} & \textbf{Dating} & \textbf{Owner} & \textbf{Tomb} & \textbf{Dating} & \textbf{Owner} \\
		\midrule
		\endfirsthead
		\multicolumn{6}{c}{\textit{Continued from previous page}} \\
		\toprule
		\multicolumn{3}{c}{\textbf{Kendall's short chronology (1999)}} & \multicolumn{3}{c}{\textbf{Török's long chronology (2009)}} \\
		\cmidrule(r){1-3} \cmidrule(l){4-6}
		\textbf{Tomb} & \textbf{Dating} & \textbf{Owner} & \textbf{Tomb} & \textbf{Dating} & \textbf{Owner} \\
		\midrule
		\endhead
		\bottomrule
		\endlastfoot
		Ku. Tum. 1 & ca. 885-835 BCE & Lord A & Ku. Tum. 1 & ca. 995-975 BCE & Prince \\
		\cmidrule(r){1-3} \cmidrule(l){4-6}
		Ku. Tum. 5 & ca. 885-835 BCE & Wife of Lord A & Ku. Tum. 4 & ca. 975-955 BCE & Prince \\
		\cmidrule(r){1-3} \cmidrule(l){4-6}
		Ku. Tum. 4 & ca. 865-825 BCE & Wife of Lord B & Ku. Tum. 5 & ca. 955-935 BCE & Prince \\
		\cmidrule(r){1-3} \cmidrule(l){4-6}
		Ku. Tum. 2 & ca. 865-825 BCE & Lord B & Ku. Tum. 2 & ca. 935-915 BCE & Prince \\
		\cmidrule(r){1-3} \cmidrule(l){4-6}
		Ku. Tum. 6 & ca. 845-815 BCE & Lord C & Ku. Tum. 6 & ca. 915-895 BCE & Prince \\
		\cmidrule(r){1-3} \cmidrule(l){4-6}
		Ku. 19 & ca. 845-815 BCE & Wife of Lord C & Ku. 19 & ca. 895-875 BCE & Prince \\
		\cmidrule(r){1-3} \cmidrule(l){4-6}
		Ku. 14 & ca. 825-805 BCE & Lord D & Ku. 14 & ca. 875-855 BCE & Prince \\
		\cmidrule(r){1-3} \cmidrule(l){4-6}
		Ku. 13 & ca. 825-805 BCE & Wife of Lord D & Ku. 13 & ca. 855-835 BCE & Prince (?) \\
		\cmidrule(r){1-3} \cmidrule(l){4-6}
		Ku. 11 & ca. 805-795 BCE & Lord E & Ku. 11 & ca. 835-815 BCE & Prince \\
		\cmidrule(r){1-3} \cmidrule(l){4-6}
		Ku. 10 & ca. 805-795 BCE & Wife of Lord E & Ku. 10 & ca. 815-795 BCE? & Unknown \\
		\cmidrule(r){1-3} \cmidrule(l){4-6}
		Ku. 9 & ca. 785 BCE & Alara & Ku. 9 & ca. 795-775 BCE & Alara \\
		\cmidrule(r){1-3} \cmidrule(l){4-6}
		Ku. 23 & ca. 785 BCE & Queen Kasaqa & Ku. 8 & ca. 775-755 BCE & Kashta \\
		\cmidrule(r){1-3} \cmidrule(l){4-6}
		Ku. 21 & ca. 760 BCE & Minor queen (?) & Ku. 17 & ca. 755-721 BCE & Piankhy \\
		\cmidrule(r){1-3} \cmidrule(l){4-6}
		Ku. 8 & ca. 760-747 BCE & Kashta & & & \\
		\cmidrule(r){1-3} \cmidrule(l){4-6}
		Ku. 7 & ca. 747-736 BCE & Queen Pebatma & & & \\
		\cmidrule(r){1-3} \cmidrule(l){4-6}
		Ku. 20 & after ca. 747 BCE & Minor queen (?) & & & \\
		\cmidrule(r){1-3} \cmidrule(l){4-6}
		Ku. 17 & ca. 716 BCE & Piankhy & & & \\
	\end{xltabular}
}

Recent excavations at the Tombos cemetery in Upper Nubia support this model, as there is little evidence of a new wave of Egyptian migration following the New Kingdom \parencite{buzonEntanglementFormationAncient2016}. In addition, a horse burial from the cemetery, dated to approximately 1005—851 BCE, is interpreted as a symbol of emerging statehood. This burial is significant because it falls between the scattered horse burials of the New Kingdom and the more ritualized horse burials at El-Kurru, associated with the pharaohs of the 25th Dynasty. As a result, this feature of Napatan culture seems to have originated in smaller communities and was later incorporated into the state \parencite{schraderColonialindigeneInteractionAncient2018}.

At present, with the evidence available, the exact dating and the origins of the first Napatan rulers remain a mystery. However, since Nubia was incorporated into the political and economic structure of the New Kingdom Egyptian administration through the authority of local elites, the Napatan kingdom could be seen as a development of the increasing power of these elites. In this sense, the “long-chronology” theory not only takes into account the authority of the indigenous elites in the aftermath of the New Kingdom collapse but also highlights Nubian agency in controlling their own sociopolitical landscape. Additionally, Nubian and Egyptian cultural contacts date to the Pre-Dynastic Period, showcasing a long history of interaction that laid the groundwork for complex cultural exchanges. Building upon this foundation of early contact and influence, we turn our attention to the broader spectrum of interactions between Nubian and Egyptian societies over the centuries.

\section{Cultural interactions between groups within Nubia and Egypt}

Historical exchanges between Nubia and Egypt have significantly influenced the development of both civilizations, with evidence of contact dating back to the 4th millennium BCE. These interactions were facilitated by extensive trade networks, through which goods and people flowed between the regions. Periods of conflict often resulted in shifts in power dynamics. As different groups and states emerged in both lands, they continued to engage with each other, leaving a lasting impact on their respective trajectories. By examining these complex interactions, we can better understand how these societies mutually shaped each other and contributed to the cultural composition of the region. We will first provide a brief overview of Nubian-Egyptian interactions through material culture and the movement of people, and then focus specifically on the New Kingdom and Napatan period, highlighting two distinct modes of interaction between these societies.

During periods of unification, the Egyptian state was interested in the Nubian region primarily to exploit its natural and human resources, and not for territorial expansion. The idea was to establish local powers under Egyptian rule in order to guarantee their economical goals \parencites[p.~94-95]{seidlmayerNubierImAgytischen2002}[p.~570]{Raue2019b}. Similarly, throughout periods of instability, e.g. First and Second Intermediate Periods, local rulers in Egypt turned towards Nubia in search of trade agreements as well as mercenaries to serve their own interests \parencite[p.~141]{naserStructuresRealitiesEgyptianNubian2013}.

The earliest Nubian society known to have had contact with Egypt, the A-Group, dates from ca. 3700—2800 BCE and occupied Lower Nubia. Other groups inhabited Upper Nubia at the same time, namely the Pre-Kerma culture and a Neolithic culture known as Akban \parencite[p.~127]{gattoAGroup4thMillennium2020}. There is no evidence of direct contact and trade between Egypt and the Pre-Kerma group \parencites[p.~148]{honeggerPreKermaCultureBeginning2020}[p.~110]{emberlingEarlyKushKingdom2022}. However, pottery evidence suggests an affinity between the A-Group and the Pre-Kerma groups \parencite[p.~222]{honeggerHolocenePrehistoryUpper2019a}. Evidence from A-Group burials, although displaying Nubian characteristics and grave goods, contained Egyptian imports \parencites{gattoNubianAGroupReassessment2006}[p.83]{williamsRelationsEgyptNubia2011}. Among them, vessels, beads, amulets, seals and combs \parencites[p.~35-36]{torokTwoWorldsFrontier2009}[p.~184-199]{royPoliticsTradeEgypt2011}.

At the same time, A-Group pottery is also attested at Egyptian Predynastic sites \parencite[p.~63-64]{gattoNubianAGroupReassessment2006}, as well as Nubian local productions such as tools and weapons made of quartz and agate in burials at the north of Aswan \parencite[p.~95]{gattoCulturalEntanglementDawn2014}. Therefore, Nubian presence in Egyptian territory is attested since the A-Group time. Settlements and cemeteries with evidence of A-Group presence were found throughout Upper Egypt, where both cultures coexisted and exhibited a blend of Egyptian and Nubian customs \parencites[p.~129]{gattoEgyptNubia5th4th2009}[p.~290-291]{meurerNubiansEgyptEarly2020}. A tumulus in Kom Ombo dated to the 4th millennium BCE has been attributed to Nubian traditions, particularly with the A-Group and Abkan cultures \parencites{gattoNubiansEgyptSurvey2005, gattoEgyptNubia5th4th2009}. An inscription from the reign of Snefru (ca. 2613—2589 BCE) refer to seven thousand Nubian prisoners taken to Egypt, while another inscription mentions 17 thousand from Lower Nubia \parencite[p.~291]{meurerNubiansEgyptEarly2020}. At Elephantine and Buhen, the contact between Nubians and Egyptians has been called “Dynastic A-Group” \parencite[p.~4]{raueWhoWasWho2008}.

Further data from Nubia include a burial at the site of Saras with a clear Egyptian influence in the form of a mud brick covering and a decorated seal found south of the First Cataract displays Egyptian iconographic symbols of royalty along with Nubian elements \parencite[p.~134, 136]{gattoAGroup4thMillennium2020}. At Cemetery L in Qustul, 25\% of pottery vessels were imported not just from Egypt, but also Palestine, including painted wares \parencite[p.~294]{Raue2019a}. It has been proposed that long-distance trade between Upper Egypt and the A-Group, with the introduction of prestige goods, played a role in the emergence of the unified Egyptian state under the 1st Dynasty (ca. 2900 BCE) \parencite[p.~49]{torokTwoWorldsFrontier2009}. Following the unification, material culture from the A-Group culture disappears, which has been attributed not only to a new Egyptian foreign policy, but also climate changes that affected the north of the continent \parencites[p.~73]{edwardsNubianArchaeologySudan2004}{gattoAGroup4thMillennium2020}.

Following the A-Group, the C-Group occupied Lower Nubia, from ca. 2500 to 1500 BCE. Evidence of long-distance trade with Egypt is found in the imported goods unearthed in their burials and graffiti left by caravan leaders serving king Pepy I \parencite[p.~166]{hafsaasCGroupPeopleLower2020}. Additionally, Egyptian tomb inscriptions and other written records provide insights into the political landscape of Nubia, naming its rulers and detailing expeditions to quarries. C-Group mercenaries were also employed in Egypt, underscoring their active involvement in Egyptian affairs \parencite[p.~96]{seidlmayerNubierImAgytischen2002}. A tableau from the Early Dynastic Period suggest the presence of Nubians in a court context in Egypt \parencite[p.~569]{Raue2019b}. Less so, are the depictions of Nubians in Egyptian domestic contexts. Two tomb paintings from the 4th Dynasty show Nubian servants as smaller than the owners of the tombs, highlighting, in Egyptian fashion, their lower social standing \parencite[p.~98-99]{seidlmayerNubierImAgytischen2002}.

During the First Intermediate Period, as Egyptian groups migrated to Lower Nubia, the influx of Egyptian imports into C-Group lands increased, including amulets, metal objects, scarabs, and pottery jars filled with grains \parencites[p.~303]{Raue2019a}[p.~171]{hafsaasCGroupPeopleLower2020}. At Elephantine, from the time of the First Intermediate Period to the 11th Dynasty, almost every major pottery collection features C-Group materials. The prevalence of these materials diminishes from the conclusion of the 12th Dynasty through to the 18th or early 19th Dynasty. In this later period, there is a significant rise in Pan-Grave and Kerma materials. Additionally, C-Group settlements have been discovered near Elephantine, along with a cemetery at Hierakonpolis, which marks the farthest northern extent of known C-Group presence in Egypt. Burial practices showcase a blend of Nubian and Egyptian traditions \parencite[p.~292]{meurerNubiansEgyptEarly2020}.

During the Middle Kingdom, the Egyptians conquered Lower Nubia, and although the C-Group initially resisted, they were eventually subjugated and collaborated with their new rulers \parencite[p.~172]{hafsaasCGroupPeopleLower2020}. This conquest allowed Egypt to acquire resources directly, reducing the need for trade and imposing stricter control on interactions between the various Nubian communities. However, the Egyptian presence in the region also benefitted the local trade \parencite[p.~138]{naserStructuresRealitiesEgyptianNubian2013}. Nubians are attested in the Egyptian court during this period, which Raue correlates to diplomatic marriages between both societies \parencite[p.~572]{Raue2019b}.

When the Egyptian centralized government collapsed during the Second Intermediate Period, evidence points to presence of the C-Group in Lower Nubia. Concurrently, other cultures emerged in Nubia, such as the Kingdom of Kerma in Upper Nubia and the Pan-Grave culture in Middle Nubia. Amidst hostilities between Egypt and Kerma during the end of the Second Intermediate Period (ca. 1555—1550 BCE), the C-Group aligned with the northern power. As a consequence, it has been suggested that the local elite adopted more Egyptian cultural markers, and even held positions of power under the Egyptian viceroy \parencites{hafsaas-tsakosKushEgyptCGroup2010}{hafsaasCGroupPeopleLower2020}.

Pan-Grave communities were a nomadic group believed to have inhabited the Eastern Desert during the 2nd millennium BCE. Archaeological evidence of their presence becomes apparent when they began interacting with Egypt and other Nubian groups, as well as their mortuary contexts \parencites[p.~81, 89]{naserNomadsNileArchaeology2012}[p.~175]{souzaNotMarginalMarginalised2022}. Their contact with Egypt is evident through the adoption of Egyptian burial customs, such as the extended burial position, and the presence of Pan-Grave pottery in Egyptian settlements and vice versa. Sites identified as Pan-Grave have been discovered in Upper Egypt, highlighting the involvement of these groups in local trade networks \parencite[p.~125-126]{souzaBlurringBoundariesMore2024}. Despite this interaction, the varying quantities of Egyptian pottery found in Pan-Grave burials suggest that their relationships with Egypt were not consistently uniform \parencite[p.~229, 237]{liszkaPanGraveMedjayIntersection2020}. Cemeteries in Upper Egypt associated with the Pan-Grave culture include Elephantine and Hierakonpolis \parencite[p.~294]{meurerNubiansEgyptEarly2020}. Additionally, further evidence from sites such as Debeira 35 in Lower Nubia and Site 4-F-74 in Upper Nubia display a mixture of burial features from the local C-Group, the mobile Pan-Grave, and migrant Egyptian populations \parencite[p.~375]{smithNubianExperienceEgyptian2020}. In Middle Egypt, their presence is associated with their employment as mercenaries for the Egyptian rulers against the Hyksos \parencite[p.~295-296]{meurerNubiansEgyptEarly2020}.

Between 2500 and 1500 BCE, the Kerma kingdom emerged in Upper Nubia's city of the same name. Trade contact with Egypt is evident by a deposit of stone vessels with Egyptian hieroglyphic writing found in the capital \parencite[p.~114]{emberlingEarlyKushKingdom2022}. This cultural exchange is further highlighted by Egyptian high-status burials containing numerous Kerma objects. Notably, the treasure of New Kingdom Queen Ahhotep featured three fly amulets crafted from gold of distinct Kushite origin, which became highly popular during the 18th Dynasty \parencite[p.~188]{Williams2020c}. Small Kerman cemeteries are known as far north as Saqqara, but no settlements have been found so far in Egyptian territory. Employment of Nubian mercenaries by the Egyptian state are suggested by Nubian material culture found in a military settlement at Deir el-Ballas \parencite[p.298-300]{meurerNubiansEgyptEarly2020}.

Conversely, Kerma tumuli tombs exhibit a blend of local and Egyptian artistic elements, including seal impressions from the Middle Kingdom and Second Intermediate Period \parencites[p.~190-191]{Williams2020c}[p.~205-206]{bonnetCitiesKermaPnubsDokki2020}. Over time, Egyptian imports become increasing common in Kerma burials and are especially high in royal tombs \parencite[p.~216-219]{bonnetEasternCemeteryKerma2020}. Elite tombs are distinguished not only by traditional Kerma features, such as bed burials and animal or human sacrifices, but also by the presence of Egyptian trade goods and locally made objects incorporating Egyptian styles \parencite[p.~118]{emberlingEarlyKushKingdom2022}. Additionally, strontium analysis of human remains from the Abu Fatima cemetery and the city of Kerma indicate a mix of local and non-local individuals during the Kerma period, with a higher proportion of non-locals at Abu Fatima (24\%) compared to Kerma (13\%) \parencite[p.~377]{schraderIntraregional87Sr86Sr2019}.

At the same time, burials also show a mixture of both Nubian cultures, Kerma and C-Group, showing contact between these groups. Objects obtained from Egypt by trade and, possibly, raids include copper daggers and Middle Kingdom statues. It is ca. 2000—1750 BCE that the first written evidence of the name “Kush” appears in Egypt. The texts depict Kerma rulers as enemies showing the hostility between both societies \parencites[p.~130-135]{emberlingPastoralStatesComparative2014}[p.~230]{honeggerHolocenePrehistoryUpper2019a}[p.~219-220]{bonnetEasternCemeteryKerma2020}[p.~130]{emberlingEarlyKushKingdom2022}. The expansion of Kerma to the north led to conflicts with Egypt, who build several fortresses in Lower Nubia during the Middle Kingdom (ca. 1985—1773 BCE) to control trade, people and the Nubian expansion \parencite[p.271]{bestockEgyptianFortressesColonization2020}. However, with the disintegration of the unified Egyptian government during the Second Intermediate Period, the fortresses were occupied by Kerma \parencites[p.~300]{meurerNubiansEgyptEarly2020}[p.~122]{emberlingEarlyKushKingdom2022}.

By the mid-second millennium, a now unified and imperialistic Egypt asserted control first over Lower Nubia and eventually to Upper Nubia, destroying the city of Kerma \parencite[p.~78]{edwardsNubianArchaeologySudan2004}. Egyptian control was extended as far as the 4th Cataract, giving the state access to Nubia’s gold mines and trade routes. By the reign of Thutmose I (1506—1493 BCE), local rulers were integrated into the Egyptian administration \parencites[p.~76]{morkotBlackPharaohsEgypts2000}[p.~99]{torokKingdomKushHandbook1997}. According to \textcite[p.~103-104]{morrisAncientEgyptianImperialism2018}, unlike the Middle Kingdom conquest, New Kingdom colonialism relied on different strategies: it aimed to transform soldiers into settlers, local elites into “Egyptianized” rulers, and the indigenous religious sphere into the cult of Amun, while reshaping local governance and economy to reflect the Egyptian model.

Egypt organized the region into Wawat for Lower Nubia and Kush for Upper Nubia, each governed by viceroys). Egypt's conquest of Nubia during the New Kingdom significantly altered Nubian-Egyptian relations, leading to an administrative control in the region and incorporation of Nubia into the Egyptian redistributive system \parencites[p.~157]{torokKingdomKushHandbook1997, torokTwoWorldsFrontier2009}{morkotConqueredConquerorOrganization2013}. Temple towns were founded throughout Nubia and provided a new venue of interactions between both cultures \parencites{kempFortifiedTownsNubia1972}{viethNewKingdomTemple2019}. According to \textcite{smithWretchedKushEthnic2003}, it is likely that these towns were inhabited by indigenous peoples, Egyptians, and those of mixed heritage through intermarriage. 

Archaeological evidence has identified several types of economic activities taking place inside the temple towns relating to the production, consumption and importing of goods, especially from Egypt \parencites{roseSesebiCeramicsChronology2017}[p.~343]{spencerBuildingNewGround2017}{cartwrightArchaeobotanicalResearchAmara2017}{taylorCoffinsDebeiraRegional2017}{budkaMaterialRemainsNew2020}. Meanwhile, evidence of Nubian imports arriving in Egypt is captured in reliefs from the tomb of the vizier Rekhmire in Thebes. These depictions feature items brought by Nubians and recorded by a scribe, including ostrich eggs and feathers, ivory, leopard skins, gold, storage jars and live animals \parencites[pl.~XVI]{daviesAncientEgyptianPaintings1936}[p.~XVIII-XX]{daviesTombRekhmiReThebes1943}[p.~301]{meurerNubiansEgyptEarly2020}.

The temples built and expanded by the Egyptians also served as economic and administrative institutions \parencite[p.~209]{torokTwoWorldsFrontier2009}. Important temples dedicated to Amun were built in cities such as Napata, Kerma, and Kawa. These temples were home to various manifestations of Amun, including a ram-headed form, which the Nubians regarded as their local version of the deity \parencite{gaboldeAmunCultIts2020}. It has been proposed that this form of Amun represents a fusion of the Egyptian god with a local ram cult from the Kerma culture \parencites[p.~77]{bonnetKermaRoyaumeNubie1990}{valbelleAmonReKerma2003}[p.~227]{torokTwoWorldsFrontier2009}. Gabolde and Török hold differing views on the extent of the Amun cult's integration into Nubian society. Török argues that the temple iconography—accessible to the general population—suggests that lower social groups were also engaged with the cult \parencite[p.~210]{torokTwoWorldsFrontier2009}. In contrast, Gabolde contends that the cult's influence was primarily confined to the elite \parencite[p.~354]{gaboldeAmunCultIts2020}. 

Archaeological data from Nubian towns provide insights into the cultural entanglement occurring there. At Amara West, local Nubian architectural traditions persist despite the settlement’s foundation by the Egyptians, and Nubian handmade pottery were used together alongside Egyptian wheelmade ones \parencite{spencerNubianArchitectureEgyptian2010, spencerMaintainingRamessideEmpire2024}. Similar patterns are seen at Sesebi and Sai Island, Egyptian-style ceramics incorporated local materials, techniques and decoration \parencites[p.~471]{roseSesebiCeramicsChronology2017}[p.~225]{budkaMaterialRemainsNew2020}. Seated female figurines discovered at Amara West reflect Egyptian influences in their form, yet their body adornments and features such as pronounced hips and inclined posture, demonstrate a connection with C-Group figurines \parencite{stevensFemaleFigurinesFolk2017}. Similarly, at Sai and Buhen, female figurines represent Egyptian styles but incised decorations manifest Nubian designs \parencite[p.~183]{budkaLifeNewKingdom2012}. This evidence illustrates that different towns experienced varying degrees of cultural entanglement and dynamics \parencite[p.~44]{spencerIntroductionHistoryHistoriography2017}. At the settlement of Askut, archaeological evidence from the New Kingdom period reveals the presence of Nubian cooking pots within domestic contexts. It has been suggested that this evidence indicates that Egyptian migrants married Nubian women, and these women maintained certain local traditions, particularly in food preparation practices \parencites{smithWretchedKushEthnic2003, smithEthnicityCulture2007}[p.~540]{vanpeltRevisingEgyptoNubianRelations2013}.

These insights into daily life and material culture also extend to burial practices, offering perspectives for examining the interplay between local and foreign traditions. At the cemeteries in Amara West, although most tombs reflect Egyptian burial traditions, one particular burial, G244, blended the Nubian tumulus with the Egyptian Egyptian underground chamber and extended body position. Additionally, grave goods included Nubian objects, differing from the contemporary Egyptian practices by focusing on daily life items rather than traditional funerary objects \parencite[p.~594-607]{binderNewKingdomTombs2017}. Isotopic analysis indicates that most individuals buried at the site, including the individual buried in G244, were local to the region \parencite[p.~44-45]{spencerMaintainingRamessideEmpire2024}. Similarly, at Sai Island, an 18th Dynasty Egyptian-style pyramid and shaft tomb belonged to Khnummose, the Overseer of Goldsmiths. Grave goods in the burial displayed Egyptian influences. However, isotopic analysis indicate that he originated from Nubia, raising the possibility that he either was Nubian in origin and later adopted Egyptian cultural markers―such as an Egyptian name and appearance―, or he descended from an Egyptian immigrant family \parencite{retzmannNewKingdomPopulation2019}.

Isotopic analysis at the Tombos cemetery further highlights burial practices across three distinct areas, showing varying proportions of local and non-local populations and their relation with certain burial practices. While Nubian tumuli contained only local individuals, Egyptian-style burials of skilled workers were mostly of non-locals. Additionally, locals were buried in both the Egyptian extended body position and the Nubian flexed position, and the same pattern is seen with burial containers; locals were not restricted to just one type, while non-locals were only interred in coffins and mats. Notably, all identified female burials at Tombos were in the Nubian flexed position, yet they were placed within Egyptian-style tombs and accompanied by Egyptian-style grave goods \parencites[p.~291-292]{buzonEntanglementFormationAncient2016}[p.~207]{smithFortifiedSettlementTombos2018}{schraderIntraregional87Sr86Sr2019} [p.~102]{smithCrossCulturalInteractionAncient2018} [p.~12-16]{buzonExploringIntersectionalIdentities2024}. The same phenomenon is observed in other sites, such as Fadrus, Qustul, Serra East and Soleb, where individuals in the Nubian flexed position are buried in Egyptian-style tombs and accompanied by Egyptian-style grave goods \parencites{schiffgiorginiSolebIINecropoles1972}{save-soderberghNewKingdomPharaonic1991}\parencite{williamsExcavationsAbuSimbel1992, williamsExcavationsSerraEast1993}. This highlights the complex interplay between cultural identity and funerary practices. Therefore, the evidence at Tombos suggests that Nubians actively engaged with both local and Egyptian burial customs, demonstrating flexibility in their funerary choices rather than rigidly adhering to a single “Egyptianized” tradition \parencite[p.~6]{buzonStrontiumIsotope87Sr2013}.

Local elite tombs also offer additional perspectives beyond material culture and human remains. Their decorated tombs provide visual representations that illustrate cultural entanglement from a different social perspective. An example is the tomb of Hekanefer in Nubia, who served the Egyptian New Kingdom administration as  “Prince of Aniba”. In his tomb, he is depicted in an Egyptianized style \parencite{simpsonHekaneferDynasticMaterial1963}. Interestingly, Hekanefer also appears in the tomb of Huy in Thebes, the Viceroy of Kush, presenting tribute to King Tutankhamun. In this depiction, he wears Egyptian dress while retaining Nubian cultural markers—ostrich feathers on his head, penannular bracelets, earrings, and animal skin. Similarly to Hekanefer’s tomb in Nubia, the local princes, Djehutyhotep and Amenemhat, were also represented as Egyptians in their own tombs in Debeira  \parencites[pl.~XXVII]{davies-gardinerTombHuyViceroy1926}[p.~269-272]{torokTwoWorldsFrontier2009}[p.~536]{vanpeltRevisingEgyptoNubianRelations2013}[p.~5]{lemosForeignObjectsLocal2020}. Scholars interpret this phenomenon as a sign of a “double or hybrid identity”, negotiated through the selective adoption and rejection of foreign practices \parencite{smithHekaneferLowerNubian2015}. On one hand, an Egyptianized depiction in a Nubian elite tomb signified higher status, prestige, and integration into the Egyptian administration. On the other, in a foreign context, Nubian elites asserted their identity through distinct cultural markers, maintaining a balance between assimilation and self-representation \parencites{torokTwoWorldsFrontier2009}{vanpeltRevisingEgyptoNubianRelations2013}.

Lemos’ analysis of grave goods in elite and non-elite cemeteries reveals significant differences in how identity and status were negotiated in New Kingdom Nubia. Elite burials, such as those of Khnummose at Sai Island and Weser in Aniba, demonstrate access to restricted Egyptian objects, such as shabtis and heart scarabs, and how locally available jars were repurposed as "canopic jars," demonstrating how imported and local objects were adapted and transformed within the burial context \parencites[p.~7]{lemosAlternativesColonizationMarginal2021}[p.~76-77]{lemosCulturalEntanglementsExperiencing2024}. Weser, in particular, possessed two heart scarabs inscribed with his name. In Egypt, heart scarabs were traditionally limited to one per burial due to their ritual connection to the individual. In Nubia, however, the presence of multiple heart scarabs likely served as a display of status, reinforcing social relations within the local community. In contrast, the material scarcity of non-elite populations constrained their access to goods, limiting their ability to negotiate identity and status through material culture. In the Fadrus cemetery, the only heart scarab was found in Tomb 511, a large two-chambered collective tomb. Lemos suggests that acquiring such a restricted object required communal effort, challenging the rigid social hierarchies imposed by colonial rule \parencites{lemosHeartScarabsOther2021, lemosCulturalEntanglementsExperiencing2024}. Furthermore, by examining patterns—or the absence thereof—in the distribution of grave goods across burials at Aniba, Sai, Soleb, and Fadrus, Lemos revealed  internal diversity within the burial assemblages. This diversity suggests that cultural entanglements were not uniform but instead shaped by local contexts. Such variations may reflect not only regional preferences but also constraints related to the consumption of specific objects by particular social groups, as exemplified by the heart scarabs mentioned earlier \parencite{lemosForeignObjectsLocal2020}.

Following the Egyptian withdrawal, archaeological evidence indicates that temple towns and cemeteries remained inhabited, suggesting that interactions between Egyptians and Nubians residents persisted even after the colonial administration disappeared. For the Napatan period, establishing a definitive chronology for cultural interactions between Nubia and Egypt remains challenging, as explained previously. The dating of the earliest burials at El-Kurru is uncertain, and at Sanam—believed to have been in use from the New Kingdom through the Napatan period—absolute dating for individual burials has not been determined \parencite{lohwasserAspekteNapatanischenGesellschaft2012}. Similarly, tombs at Meroe West and South, dated solely by Reisner during his early 20th-century excavations, present significant ambiguities. A considerable number, particularly those without grave goods, remain “undated”. As a result, a precise chronology for analyzing material culture—from the beginning of the period, through the Egyptian conquest by the Napatan rulers, their eventual withdrawal, and the transition to the Meroitic period—cannot be established.

Regardless, Egyptian-style objects, as well as Egyptian-style burial practices, are present in Nubian cemeteries, attesting to the ongoing interactions between Egyptians and Nubians attested though material culture.	The cemetery of Tombos, for example, remained in use during the Napatan period. Egyptian-style burials, featuring pyramids and underground chambers, continued alongside Nubian-style tumuli. Most burials followed the Egyptian extended position. Grave goods included Egyptian-style objects, such as an inscribed heart scarab found near a pyramid and a scarab bearing the name of Piankhy \parencite[p.~293-294]{buzonEntanglementFormationAncient2016}. Traditional Nubian objects include Red Sea shell beads, ivory bracelets and quartz crystal pendants, that resemble those dated to the Kerma Period. Particularly interesting objects include a headrest similar to those found in Kerma burials—these were already adaptations of typical Egyptian headrests—and a heart scarab that showcases unique artistic practices. These practices reference Egyptian styles but do not simply imitate them. The pottery from burials showcases a fusion of Nubian and Egyptian influences, combining decorative motifs and manufacturing techniques from both cultures \parencites[p.~6-9]{smithColonialEntanglementsEgyptianization2014}[p.~637-642]{buzonTumuliTombosInnovation2023}.

The tumuli typically held one individual, a practice linked to burial traditions at Kerma. In these burials, Egyptian-style coffins were often paired with Nubian burial beds. One woman was buried in the Nubian flexed position on a burial bed. Isotopic analysis revealed that none of the Napatan samples had Egyptian strontium values, indicating that Egyptian immigration had ceased by this time. However, morphological and isotopic evidence suggests that Nubians from other regions were integrated into the Tombos community. Researchers propose that this population may represent descendants of the mixed Nubian-Egyptian community from the New Kingdom \parencites{schraderIlluminatingNubianDark2014}[p.~294-295]{buzonEntanglementFormationAncient2016}{buzonTumuliTombosInnovation2023}{buzonExploringIntersectionalIdentities2024}{smithCrossCulturalInteractionAncient2018}.

Similar to Tombos, the burials from the Napatan period at Sedeinga feature both tumuli and pyramids, with the latter being reused during the Meroitic period \parencites{berger-elnaggarContributionSedeingaLhistoire2008}{rillyCloserAncestorsExcavations2018}. Recent excavations have revealed a particularly interesting burial. The undisturbed grave II TT 262 belonged to a very young individual less than two years old, who was interred in a lateral niche grave topped by a tumulus superstructure. Interestingly, this grave was related to two pyramids, as it was constructed behind them. This placement highlights the blend of local tumulus traditions with the architectural style of Egyptian pyramids. The grave goods included an Egyptian-style amulet depicting the god Shu and copper-alloy anklets. Napatan period ceramics from the necropolis comprise both Egyptian imports and local productions \parencites{rillySedeinga2012Season2013, rillyCloserAncestorsExcavations2018, rillyCollectiveGravesBastatues2020}.

The transition cemeteries of Amara West—Cemeteries C and D—feature tombs spanning from the New Kingdom to the early Napatan period. This site, like the other cemeteries in the region, exhibits a similar blend of Nubian and Egyptian burial traditions. Egyptian chamber tombs, for instance, incorporated both Egyptian elements, such as headrests and coffins, and Nubian cultural features, including burial beds and flexed body positions. Similarly, Nubian tumuli tombs often contained coffins and extended burials. Both types of burials contained Egyptian-style amulets \parencites{spencerAmaraWestII2002}{binder10th9thCenturyBC2011}. At the site of Kirbekan in the region of the 4th Cataract, pottery assemblages from Napatan tombs also attested to a blend of Nubian and Egyptian traditions \parencite{budkaDocumentationExcavationDome2007}.

During the Kushite 25th Dynasty in Egypt, evidence of Nubian burials in the region is found in Abydos and Thebes \parencites{budkaKuschitenAbydosEinige2012}{budkaNubians1stMillennium2019}. At Abydos, five monuments attest to the presence of Nubians in Egypt at this period. This includes a stela depicting the Kushite prince Ptahmaakheru, who likely belonged to the royal family. Other monuments feature another member of the royal family and an official or member of the royal household \parencite{leahyKushitesAbydosRoyal2014}. These can be identified by their non-Egyptian name, costume and physical depiction \parencite[p.~704-705]{budkaNubians1stMillennium2019}. In Thebes, a cluster of Kushite tombs, as well as intrusive Kushite burials in Tomb TT 99 of Sennefer, was discovered. Notably, this included the coffin of Niw, daughter of Padiamun, Priest of Amun in the Nubian city of Kawa. Her depiction reveals a blend of Egyptian and Nubian features, such as an Egyptian garment paired with a short Kushite hairstyle. Additionally, bandages belonging to other Kushite officials were found among the intrusive burials \parencite[p.~504-505]{budkaKushiteTombGroups2010}.

Another tomb, featuring a mud brick superstructure, contained the remains of several individuals of Kushite ancestry. Although much of the burial assemblage reflects Egyptian characteristics, certain distinct Kushite elements were preserved. These include personal names, depictions, Nubian objects, and Egyptian grave goods with unique Nubian attributes, all of which highlight the non-Egyptian origin of these individuals. Depictions of Kushites in Egypt can also exhibit Egyptian, Nubian, or a blend of both features. In some cases, it is only through their father's name that we recognize them as Nubian, suggesting that these individuals had the flexibility to adapt or change their identity as they wished \parencites{budkaKushiteTombGroups2010, budkaNubians1stMillennium2019, budkaNubiansEgypt25th2020}.

The apparent absence of Kushite artifacts in Egypt could be due to either a lack of resources to personalize Egyptian objects or a lack of interest in doing so, which contributed to a uniformity in the archaeological record \parencite[p.~494]{howleyKushitesEgypt6642020}. While archaeological evidence for Nubians—especially from lower social groups—is scarce, it is likely that they settled in Egypt during the 25th Dynasty. Additionally, some may have married Egyptians, leading to the formation of intercultural families similar to those found in Nubian towns \parencite[p.~477]{budkaNubiansEgypt25th2020}.

25th Dynasty royal burials offer further insights into cultural practices and interactions during the period. A re-examination of shabtis uncovered a particular type featuring a basket on its head, discovered exclusively in three queens’ tombs at El-Kurru (Ku. 62, 71 and 72) and an elite New Kingdom tomb in Thebes (TT 99) \parencite{mussoKushiteShabtisBasket2011}. Furthermore, imported marl pottery jars from El-Kurru indicate trade connections with Upper Egypt, which apparently ceased during the later half of the 7th century BCE \parencite{heidornBostonMuseumFine2018}. Additionally, pot-breaking ceremonies observed in tomb enclosures at El-Kurru tombs mirror those performed by Egyptian priests at Abydos \parencite{budkaEgyptianImpactPotBreaking2014}. There rituals, already present in Nubia since the New Kingdom, underscore a cultural continuity and exchange \parencites[p.~154, 165-166]{smithWretchedKushEthnic2003}[p.~120]{gasperiniAmaraWestCemeteries2023}.

After the Kushite kings withdrew from Egypt around 664 BCE, textual and archaeological sources confirm the ongoing presence of Nubian officials who retained their positions during the 26th Dynasty. Notable figures included the God's Wife of Amun; the High Priest of Amun and grandson of Shabaqo; the wife of the Mayor of Thebes, who was also the granddaughter of Piankhy; and Peksater, a wife of Piankhy. Furthermore, Egyptians began incorporating Nubian features into their burials, such as the bead net and offering tables, with the latter having been absent in Egypt from the 21st to the 24th Dynasties. The presence of Egyptian-style objects in the Nuri royal necropolis in Nubia also attests to the continued trade and movement of specialized people between the two regions \parencites[p.~514]{budkaKushiteTombGroups2010}[p.~494-498]{howleyKushitesEgypt6642020}.

During the 6th century BCE, relations between Nubia and Egypt were strained, likely due to the invasion of the 26th Dynasty king Psamtek II into Lower Nubia, which eventually became part of Egypt. During this time, the use of Egyptian writing became less common. Trade relations between Nubia and Egypt were reestablished around 570—526 BCE. Following the Persian conquest of Egypt, a text from the reign of Darius I (ca. 522—486 BCE) records a shipment of ivory from Kush. Additionally, it is likely that Nubians served in Xerxes I’s army. During the reign of Harsiyotef (ca. 404-369 BCE)—one of the last kings of the Napatan period—, a stela inscribed with Egyptian hieroglyphics documents military campaigns in Lower Nubia. The text details battles against non-Egyptian groups, highlighting the region's distinctly intercultural nature even after its integration into Egypt \parencites[p.~344-345, 344-345, 376-377]{torokKingdomKushHandbook1997}[p.~361-370]{torokTwoWorldsFrontier2009}[p.~498-499]{howleyKushitesEgypt6642020}.

Beyond the royal pyramids, non-elite tombs from the later Napatan period are absent from the archaeological record. This absence can be attributed to several factors. Only Reisner’s work at the Meroe cemeteries included non-royal burials dated to the Late Napatan phases, but these were exclusively elite populations. Therefore, the excavation methods of the time may have been inadequate to identify and date the simpler graves associated with non-elites. Additionally, the lack of datable material culture in poorer tombs might have caused them to remain undated os misdated to earlier phases. Other possibilities include the existence of undiscovered later Napatan cemeteries and the reuse of tombs over time may have hindered their identification.

\section{Conclusion}

In this chapter, I explored the extensive history of cultural interactions between Nubia and Egypt, which shaped both cultures for millennia. This historical background is crucial as it demonstrates that the presence of Egyptian-style material culture in Napatan burials is not just an isolated occurrence bur rather a continuation of a long-standing process of cultural exchange. In particular, the position of Nubia in relation to Egypt during the New Kingdom period significantly influenced the development of material culture of the Nubian populations in subsequent years. This complex historical context is essential for its influence in material culture assemblages of the Napatan period. Especially how different people selectively adopted and adapted certain cultural markers, reflecting their responses to shifting dynamics with Egypt.

The historical background also reveals the intricate and evolving relationships between Kush and Egypt over time during the Napatan period. These ranged from cycles of Kushite control over Upper Egypt, periods of trade and conflict with the Assyrians, and to religious and political relations with the Amun cult in Thebes. These shifts set a stage for how amulet use might have varied across different social groups, depending on the political and cultural ties with Egypt. By examining these dynamics through a chronological and regional lens, we can better understand how amulet use shifted in response to these changing contacts with Egypt.
